\chapter{Construire un éditeur de texte, puis de code}

Jusqu'ici, l'utilisateur \emph{regardait}. Ici il \textbf{modifie}. Tout change :
il faut un curseur, une sélection, un défilement, une annulation --- et il faut
décider comment le texte est rangé en mémoire, décision qui gouvernera tout le
reste.

Nous construisons l'éditeur en douze étapes. Aux étapes 1 à 9 vous avez un
éditeur de texte utilisable ; aux étapes 10 à 12 il devient un éditeur de code.

La référence est \texttt{NkCodeEditor.h} du dépôt : \textbf{5\,012 lignes}, plus
\texttt{NkSyntax.h} (1\,422) pour la coloration.

\begin{nkretenir}
\textbf{Les étapes 1 et 2 du chapitre 21 ne changent pas.} Ajoutez seulement, aux
dépendances du \texttt{.jenga}, \texttt{NKGui} et \texttt{NKFileSystem} --- ils y
sont déjà --- et créez \texttt{src/Editeur/}.
\end{nkretenir}

\section{Étape 1 --- décider comment ranger le texte}

C'est la première décision, et elle est irréversible en pratique.

\begin{center}
\begin{tabular}{@{}lll@{}}
\toprule
\textbf{Représentation} & \textbf{Bon marché} & \textbf{Cher} \\
\midrule
une seule chaîne & lire tout d'un coup & insérer au milieu \\
vecteur de lignes & éditer une ligne, dessiner & insérer une ligne au début \\
tampon à trou (\emph{gap}) & frappe au clavier & sauts lointains \\
\bottomrule
\end{tabular}
\end{center}

\begin{nkretenir}
\textbf{Choisissez d'après l'opération la plus fréquente, pas d'après la plus
élégante.} Un éditeur passe son temps à \emph{dessiner des lignes} et à
\emph{éditer celle du curseur} ; il n'insère presque jamais un million de lignes
en tête.

NKCode prend donc le \textbf{vecteur de lignes}. L'opération chère n'arrive
jamais.
\end{nkretenir}

\begin{nkcode}{Applications/NKCode/src/NKCode/Editor/NkCodeEditor.h}
// Une ligne = caracteres bruts (PAS de '\0').
using NkCodeLine = NkVector<char>;
\end{nkcode}

\begin{nkpiege}
\textbf{« PAS de \texttt{'\textbackslash 0'} »} : le commentaire d'origine le dit
en majuscules. Une ligne est une plage $[\texttt{begin}, \texttt{end})$, pas une
chaîne C.

Passer \texttt{line.Data()} à une fonction qui attend une chaîne terminée lirait
au-delà de la ligne, jusqu'au premier zéro rencontré dans le tas --- donc parfois
rien d'anormal, et parfois un plantage. Le pire des deux mondes.
\end{nkpiege}

\section{Étape 2 --- le document, et où vit l'état}

\begin{nkcode}{Applications/NKCode/src/NKCode/Editor/NkCodeEditor.h}
struct NkCodeDoc {
        NkVector<NkCodeLine> lines;
        int32 curLine = 0, curCol = 0;   // curseur (col = index caractere)
        int32 selLine = 0, selCol = 0;   // ancre de selection (== curseur si vide)
        // ... defilement, etat d'annulation
};
\end{nkcode}

L'en-tête explique \emph{pourquoi} le curseur est là et pas ailleurs :

\begin{nkcode}{Applications/NKCode/src/NKCode/Editor/NkCodeEditor.h}
// etat curseur/selection/scroll EMBARQUE dans le document -> survit au
// changement d'onglet ET au re-dock (etat hors-frame).
\end{nkcode}

\begin{nkretenir}
\textbf{L'interface est en mode immédiat : elle ne garde rien d'une image à
l'autre.} Un état rangé dans le widget disparaîtrait au premier changement
d'onglet --- vous reviendriez sur votre fichier, curseur au début, défilement en
haut.

\textbf{Ce qui doit survivre à la disparition du widget vit dans le document.}
C'est la règle générale du mode immédiat, et l'éditeur en est le cas le plus
visible.
\end{nkretenir}

\begin{nkretenir}
\textbf{La sélection est une ANCRE plus le curseur, pas un couple début/fin.}
L'ancre ne bouge pas pendant \texttt{Maj+flèches} ; le curseur, oui. La sélection
est l'intervalle entre les deux, dans l'ordre où il se trouve.

Stocker « début » et « fin » obligerait à savoir lequel des deux bouge --- et
c'est exactement l'information que l'ancre porte gratuitement.
\end{nkretenir}

\section{Étape 3 --- charger et enregistrer}

\begin{nkcode}{La forme}
bool Charger(const NkString &chemin, NkCodeDoc &doc);   // decoupe sur \n
bool Enregistrer(const NkString &chemin, const NkCodeDoc &doc);
\end{nkcode}

\begin{nkpiege}
\textbf{Les fins de ligne : \texttt{\textbackslash n} ou
\texttt{\textbackslash r\textbackslash n}.} Un fichier Windows en porte deux
octets, un fichier Unix un seul.

Découper naïvement sur \texttt{\textbackslash n} laisse un
\texttt{\textbackslash r} \textbf{à la fin de chaque ligne}. Il est invisible à
l'affichage --- et il ressort à l'enregistrement, doublé, puis quadruplé au
suivant.

Décidez à l'ouverture, retenez ce que le fichier employait, et réécrivez la
même chose. \emph{Un éditeur qui change les fins de ligne rend un diff
illisible.}
\end{nkpiege}

\section{Étape 4 --- ne dessiner que ce qui se voit}

\begin{nkcode}{Le calcul, avant toute boucle de dessin}
const float32 hauteurLigne = mPolice.LineHeight();
const int32 premiere = (int32)(doc.scrollY / hauteurLigne);
const int32 combien  = (int32)(zone.h / hauteurLigne) + 2;   // +2 : les moities
const int32 derniere = NkMin(premiere + combien, (int32)doc.lines.Size());

for (int32 i = premiere; i < derniere; ++i) { DessinerLigne(dl, doc.lines[i], ...); }
\end{nkcode}

\begin{nkretenir}
\textbf{Le coût du dessin devient indépendant de la taille du fichier} :
cinquante lignes, qu'il y en ait cent ou cent mille.

Le \texttt{+2} n'est pas de la superstition : la première ligne visible est
souvent coupée en haut, la dernière en bas. Sans lui, une bande vide apparaît au
bord pendant le défilement.
\end{nkretenir}

\begin{nkretenir}
L'en-tête de NKCode assume une approximation et l'écrit : « largeur H approximée
sur les lignes visibles ». La largeur de défilement horizontal n'est calculée que
sur ce qui est à l'écran, parce que mesurer cent mille lignes à chaque image
coûterait plus que le bénéfice.

\textbf{Une approximation écrite est une décision ; une approximation tue est un
défaut en attente.}
\end{nkretenir}

\section{Étape 5 --- le curseur, et la conversion des coordonnées}

\begin{nkcode}{Les deux conversions, et il faut les deux}
// ecran -> document : ou l'utilisateur a-t-il clique ?
int32 ligne = (int32)((p.y - zone.y + doc.scrollY) / hauteurLigne);
int32 col   = ColonneLaPlusProche(doc.lines[ligne], p.x - zone.x + doc.scrollX);

// document -> ecran : ou dessiner le caret ?
float32 x = zone.x - doc.scrollX + LargeurTexte(doc.lines[curLine], 0, curCol);
float32 y = zone.y - doc.scrollY + curLine * hauteurLigne;
\end{nkcode}

\begin{nkpiege}
\textbf{Les deux conversions doivent être inverses l'une de l'autre, et rien ne
le vérifie.} Si le dessin soustrait le défilement et que le clic ne l'ajoute pas,
la sélection est juste en haut du fichier et fausse plus bas --- donc juste
pendant tout le développement.

C'est le même piège que le liseré de sélection peint en espace document au lieu
d'espace écran : deux chemins qui convertissent entre les mêmes espaces
\textbf{doivent partager leur fonction de conversion}, pas la recopier.
\end{nkpiege}

\begin{nkretenir}
\textbf{\texttt{ColonneLaPlusProche} arrondit au caractère le plus proche, pas au
caractère survolé.} Cliquer sur la moitié droite d'une lettre place le curseur
\emph{après} elle --- c'est ce que fait tout éditeur, et ne pas le faire donne
une impression de décalage d'un caractère.
\end{nkretenir}

\section{Étape 6 --- la saisie}

\begin{center}
\begin{tabular}{@{}ll@{}}
\toprule
\textbf{Entrée} & \textbf{Effet sur le document} \\
\midrule
caractère & insérer à \texttt{(curLine, curCol)}, \texttt{++curCol} \\
\texttt{Entrée} & couper la ligne en deux, descendre \\
\texttt{Retour arr.} & si \texttt{curCol > 0} effacer ; sinon \emph{fusionner} \\
\texttt{Suppr} & symétrique, vers l'avant \\
flèches, Home, End & déplacer, et poser l'ancre \emph{sauf} avec Maj \\
\bottomrule
\end{tabular}
\end{center}

\begin{nkretenir}
\textbf{Les deux cas de bord sont la fusion de lignes.} \texttt{Retour arrière}
en colonne 0 ne supprime pas un caractère : il colle la ligne courante à la
précédente et place le curseur à la jointure. \texttt{Suppr} en fin de ligne fait
le symétrique.

Ce sont les deux seules opérations qui changent le \emph{nombre} de lignes en
dehors d'\texttt{Entrée}. Elles sont oubliées une fois sur deux, et le symptôme
est « la touche ne fait rien en début de ligne ».
\end{nkretenir}

\begin{nkpiege}
\textbf{Toute saisie avec sélection non vide commence par supprimer la
sélection.} Taper une lettre alors qu'un bloc est sélectionné doit le remplacer.

Oublier cette règle donne un éditeur où la sélection ne sert qu'à copier --- et
le défaut ne se voit pas tant qu'on ne teste que la frappe simple.
\end{nkpiege}

\section{Étape 7 --- le défilement automatique}

\begin{nkretenir}
\textbf{Après toute opération qui bouge le curseur, ramenez-le dans la vue.}
Sinon la flèche « bas » sort de l'écran et l'utilisateur tape à l'aveugle.

La règle est simple et il faut les deux moitiés : si le curseur est au-dessus de
la vue, faire remonter le défilement à sa ligne ; s'il est en dessous, le
descendre jusqu'à ce que sa ligne soit la dernière visible.
\end{nkretenir}

\section{Étape 8 --- l'annulation}

\begin{nkcode}{Applications/NKCode/src/NKCode/Editor/NkCodeEditor.h}
struct UndoState { /* ... */ };
\end{nkcode}

\begin{nkretenir}
\textbf{Deux stratégies, et la plus simple est souvent la bonne.}

\emph{Instantané} : une copie du document avant chaque groupe de modifications.
Simple, sûr, coûteux en mémoire --- mais un fichier source fait quelques dizaines
de kilo-octets, donc cent états tiennent dans quelques méga-octets.

\emph{Journal d'opérations} : « inséré \texttt{x} en (ligne, colonne) », et on
sait le défaire. Économe, et bien plus difficile à rendre juste : chaque
opération a besoin de son inverse exact, et un seul inverse faux corrompt tout
l'historique sans prévenir.

\textbf{Commencez par l'instantané.} Ne passez au journal que si une
\emph{mesure} --- pas une intuition --- montre que la mémoire pose problème.
\end{nkretenir}

\begin{nkpiege}
\textbf{Un groupe d'annulation n'est pas une frappe.} Si chaque caractère tapé
crée un état, l'utilisateur devra faire trente annulations pour effacer un mot.

Le regroupement se fait sur un \emph{critère} : une pause de frappe, un
changement de ligne, un changement de nature d'opération (saisie puis
effacement). Ce critère est un choix de conception, pas un détail --- et c'est
celui que les utilisateurs remarquent en premier.
\end{nkpiege}

\begin{nkretenir}
\textbf{L'instantané doit contenir le curseur.} Annuler et retrouver le curseur
ailleurs est désorientant : l'utilisateur veut revenir \emph{à l'endroit} où il
était, pas seulement au texte.
\end{nkretenir}

\section{Étape 9 --- la gouttière}

Numéros de ligne à gauche, alignés à droite dans une colonne dont la largeur
dépend du nombre de chiffres du fichier.

\begin{nkretenir}
\textbf{Calculez la largeur sur le plus grand numéro, pas sur une constante.}
C'est exactement le piège de la table des matières de ce livre : une boîte
dimensionnée pour trois chiffres écrase le texte à partir de la ligne 1000.

Et le défaut n'apparaît que sur les gros fichiers --- donc jamais pendant que
vous écrivez l'éditeur.
\end{nkretenir}

\noindent \textbf{Vous avez un éditeur de texte utilisable.} Les trois étapes
suivantes en font un éditeur de code.

\section{Étape 10 --- décrire un langage plutôt que le programmer}

\begin{nkcode}{Applications/NKCode/src/NKCode/Editor/NkSyntax.h}
struct NkLangSpec {
        NkString name;
        NkVector<NkString> exts;         ///< ".css", ".scss"... (minuscules)
        NkString lineComment;            ///< "//", "#", "--", ";" ("" = aucun)
        NkString blockOpen, blockClose;  ///< "/*","*/" ou "<!--","-->"
        NkString strChars;               ///< delimiteurs de chaines
        NkVector<NkString> keywords;     ///< tries
        NkVector<NkString> types;        ///< tries
};
\end{nkcode}

\begin{nkretenir}
\textbf{Un langage n'est pas du code : c'est une donnée.} Ajouter la coloration
du Lua, du Rust ou du CSS ne demande pas d'écrire un colorateur --- seulement de
remplir un \texttt{NkLangSpec}.

Un colorateur écrit langage par langage aurait donné dix fonctions qui divergent ;
une description en donne dix \emph{données} qui ne peuvent pas diverger, puisqu'un
seul code les lit.

\textbf{C'est le test de votre étape 10} : si ajouter un second langage vous fait
écrire une ligne de code, elle n'est pas finie.
\end{nkretenir}

\begin{nkretenir}
\texttt{keywords} et \texttt{types} sont \textbf{triés} : la recherche est
dichotomique. Sur un langage à 90 mots-clés, testés pour chaque mot de chaque
ligne visible, la différence entre parcours linéaire et dichotomie se voit à
l'\oe il pendant le défilement.
\end{nkretenir}

\section{Étape 11 --- tokeniser une ligne, sans allouer}

\begin{nkcode}{Applications/NKCode/src/NKCode/Editor/NkSyntax.h}
// Tokenise UNE ligne et emet des plages colorees [begin,end) via un callback
// Sans allocation : gap-filling (les zones non colorees -> couleur texte)
\end{nkcode}

\begin{nkretenir}
\textbf{Deux décisions de performance dans deux lignes de commentaire.}

\emph{Une ligne à la fois} : la coloration se fait au moment de dessiner, sur les
lignes visibles seulement. Rien n'est stocké, donc rien ne peut être périmé après
une modification --- pas de cache à invalider, pas de cache qui ment.

\emph{Sans allocation, par remplissage de trous} : au lieu de produire la liste
de toutes les plages --- y compris celles qui gardent la couleur par défaut ---
le tokeniseur n'émet que les plages \emph{colorées}, et ce qui reste entre elles
prend la couleur du texte.
\end{nkretenir}

\section{Étape 12 --- l'état d'entrée de ligne}

\begin{nkpiege}
\textbf{Un commentaire de bloc ne tient pas dans une ligne.} \texttt{/*} ouvert
ligne 10 et fermé ligne 40 : les trente lignes du milieu sont dans le
commentaire, et \emph{rien dans leur texte ne le dit}.

Un tokeniseur strictement par ligne ne peut pas le savoir : il lui faut un
\textbf{état d'entrée de ligne} --- « suis-je déjà dans un commentaire ? » ---
calculé depuis le début du fichier, ou depuis un point de reprise connu.

C'est la seule vraie difficulté de la coloration syntaxique, et c'est celle qu'on
découvre en dernier, quand un fichier entier devient vert.
\end{nkpiege}

\begin{nkretenir}
\textbf{Le remède pratique} : garder, pour chaque ligne, un octet d'état de
sortie. Le calculer une fois au chargement, et le mettre à jour à partir de la
ligne modifiée jusqu'à ce que l'état redevienne identique à l'ancien --- ce qui
arrive presque toujours en une ou deux lignes.

Recalculer tout le fichier à chaque frappe marche aussi, et se voit dès dix mille
lignes.
\end{nkretenir}

\begin{nkbilan}
Un éditeur repose sur douze étapes, et les quatre premières décident du reste :
la représentation du texte (choisie d'après l'opération \emph{la plus fréquente}),
le lieu de l'état d'édition (dans le document, jamais dans le widget), les fins
de ligne (retenues et restituées), et le dessin restreint aux lignes visibles.
Viennent ensuite les deux conversions écran/document, qui \textbf{doivent être
inverses} et donc partager leur code ; la saisie, dont les deux cas de bord sont
les fusions de lignes ; l'annulation, par instantané avant tout journal. Passer
au code n'ajoute qu'une chose : la coloration --- et le bon réflexe est de
\emph{décrire} chaque langage par une donnée. La seule vraie difficulté qui reste
est l'état d'entrée de ligne.
\end{nkbilan}

\begin{nkquestions}
\begin{enumerate}
\item Citez trois représentations d'un texte éditable et l'opération que chacune
      rend chère.
\item Pourquoi l'état du curseur ne peut-il pas vivre dans le widget ?
\item Pourquoi une ligne ne se termine-t-elle pas par
      \texttt{'\textbackslash 0'} ?
\item Que doivent faire \texttt{Retour arrière} en colonne 0 et \texttt{Suppr}
      en fin de ligne ?
\item Pourquoi les deux conversions écran/document doivent-elles partager leur
      code ?
\item Pourquoi un tokeniseur par ligne a-t-il besoin d'un état d'entrée de
      ligne ?
\end{enumerate}
\end{nkquestions}

\begin{nkpratique}
Construisez votre éditeur en suivant les douze étapes, sans en sauter une.

\textbf{Étape A (1 à 9) --- le texte.} Ouvrir un fichier, l'afficher avec ses
numéros de ligne, un curseur, les flèches, \texttt{Home}, \texttt{End}, la
saisie, \texttt{Entrée}, \texttt{Retour arrière}, \texttt{Suppr}, la sélection à
la souris et au clavier, le défilement à la molette, l'annulation, et
l'enregistrement.

\textbf{Étape B (10 à 12) --- le code.} Un \texttt{NkLangSpec} pour un langage,
et la coloration des mots-clés, types, chaînes, nombres, commentaires de ligne.

\textbf{Le critère de réussite de l'étape B est binaire} : ajoutez un
\emph{second} langage sans écrire une seule ligne de code de coloration. Si vous
n'y arrivez pas, votre étape 10 n'était pas assez « décrite ».
\end{nkpratique}

\begin{nkdemo}
Prouvez que votre éditeur ne dessine que les lignes visibles.

Ajoutez un compteur de lignes réellement dessinées, affiché à l'écran. Ouvrez un
fichier de dix lignes, puis un de cent mille lignes générées. Le compteur doit
être \emph{le même} dans les deux cas.

Mesurez ensuite le temps d'une image avec \texttt{NkChrono}, sur les deux
fichiers, dix fois chacun.

\textbf{Ce que la démonstration doit établir} : que l'écart entre les deux temps
reste sous le bruit de mesure --- et il vous faut donc \emph{d'abord} mesurer ce
bruit, en comparant deux séries sur le \emph{même} fichier. Sans ce témoin, vous
ne pouvez pas dire si un écart de 3\,\% est un effet ou du hasard.
\end{nkdemo}

\begin{nklogique}
Un utilisateur tape \texttt{Bonjour}, appuie sur \texttt{Entrée}, tape
\texttt{monde}, puis fait trois annulations.

\begin{enumerate}
\item Décrivez l'état obtenu selon trois politiques de regroupement : par
      caractère, par mot, par ligne.
\item Laquelle vous paraît la meilleure ? Défendez-la, puis donnez le cas où elle
      surprend l'utilisateur.
\item Le document est ensuite \emph{rechargé depuis le disque} par une action
      extérieure. L'historique d'annulation doit-il survivre ? Argumentez dans
      les deux sens, tranchez, et dites ce que votre choix rend impossible.
\end{enumerate}
\end{nklogique}
